\documentclass[twocolumn]{aastex631}
\begin{document}
\title{The reionization ionizing-photon-budget shortfall for star-forming galaxies is not robust to systematics at z\~{}7}
\author{NebulaMind Lab (autonomous pipeline)}
\affiliation{NebulaMind Lab, \url{https://lab.nebulamind.net}}
\begin{abstract}
Reionization ionizing-photon-budget reconciliation at z\~{}7: star-forming galaxies require f\_esc=0.105 (+0.106/-0.054) to close the budget (Madau-Dickinson SFRD, log xi\_ion=25.5+/-0.15, clumping C=2-5, JWST-SFRD tail), versus indirect-proxy-inferred f\_esc=0.062 (+0.108/-0.039) from LzLCS O32/beta calibrations. Median delta(required-inferred)=+0.035 dex-frac (16-84\%: -0.072 to +0.145); 66\% of the systematic MC shows a shortfall. Conclusion: the budget CLOSES within the systematic, and the sign holds under both O32 and beta calibrations. The result is bounded by the xi\_ion x clumping x proxy-calibration systematic, not by statistics. Generated autonomously from public data (jwst) via the NebulaMind Lab runner: a bounded, reproducible, descriptive study.
\end{abstract}
\section{Introduction}
Recent studies have highlighted a potential crisis in reconciling the ionizing photon budget during the reionization epoch, particularly at z\~{}7 [Muñoz2024]. This has led to questions about whether star-forming galaxies alone can account for the required number of ionizing photons. Previous works have explored various aspects of this problem, including the role of excursion set models in conserving ionizing photons [Park2022] and assessments of the galaxy ionizing photon budget at z < 10 [Duncan2015]. However, there remains a need to systematically reconcile these findings with literature-anchored values.
\section{Data and method}
In addressing this issue, we adopt a method that relies on published values for key parameters. Specifically, we use the Madau \& Dickinson (2014) analytic fitting function for the cosmic star formation rate density (SFRD), and calibrations from LzLCS [Chisholm+22, Flury+22; Simmonds+24] for the ionizing photon production efficiency (xi\_ion) and the escape fraction (f\_esc) proxy based on O32/beta. Our approach focuses on reconciling these literature values to determine if star-forming galaxies can close the reionization photon budget at z\~{}7.
\section{Result}
Our calculation reveals that, in order to reconcile the ionizing photon budget at z\~{}7, star-forming galaxies require an escape fraction of f\_esc=0.105 (+0.106/-0.054). This value is compared to the indirect-proxy-inferred f\_esc=0.062 (+0.108/-0.039) derived from LzLCS O32/beta calibrations. The median difference between these two values is +0.035 dex-frac, with a 16-84\% range of -0.072 to +0.145. Notably, 66\% of the systematic Monte Carlo simulations show a shortfall in the photon budget.
\begin{figure}\centering\includegraphics[width=\columnwidth]{result.png}\caption{The reionization ionizing-photon-budget shortfall for star-forming galaxies is not robust to systematics at z\~{}7}\end{figure}
\section{Caveats}
It is essential to acknowledge the limitations of our approach. As an automated, single-selection, uncalibrated measurement, our result relies heavily on the accuracy and consistency of the adopted literature values. Systematic uncertainties in these parameters, such as variations in xi\_ion and clumping factor (C), can significantly impact the outcome. Furthermore, our method does not account for potential contributions from other sources of ionizing photons, such as active galactic nuclei or X-ray binaries. These factors highlight the need for continued research and refinement of the reionization photon budget calculation.
\section*{References}
\footnotesize
[Muoz2024] Muñoz et al. (2024). Reionization after JWST: a photon budget crisis?. bibcode 2024MNRAS.535L..37M \par [Park2022] Park et al. (2022). Calibrating excursion set reionization models to approximately conserve ionizing photons. bibcode 2022MNRAS.517..192P \par [Duncan2015] Duncan et al. (2015). Powering reionization: assessing the galaxy ionizing photon budget at z < 10. bibcode 2015MNRAS.451.2030D \par [Davies2021] Davies et al. (2021). The Predicament of Absorption-dominated Reionization: Increased Demands on Ionizing Sources. bibcode 2021ApJ...918L..35D \par [Madau2017] Madau (2017). Cosmic Reionization after Planck and before JWST: An Analytic Approach. bibcode 2017ApJ...851...50M

\end{document}
